\chapter{Stratégies de flux de change}\label{ch:s2:fx-flow-strategies}

À seize heures, heure de Londres, les taux de référence des devises sont fixés sur les transactions de quelques minutes. Les fonds qui doivent négocier au fixing
créent des flux que d'autres peuvent anticiper, à partir d'informations publiques sur ce que ces fonds détiennent. Melvin et Prins ont établi qu'un mois au cours duquel les marchés d'actions
progressent prédit une dépréciation de la devise avant le fixing de fin de mois, car les investisseurs internationaux en actions rééquilibrent leurs couvertures de change. Sur des données publiques
de 1999 à 2026, lorsque le S\&P 500 avait progressé au cours d'un mois, le dollar avait tendance à baisser pendant les deux derniers jours du mois. Un trade qui le vendait alors
rapportait 11.4~points de base par mois ($t$ = 3.4). Sur le marché synthétique de ce chapitre, prendre position avant un flux de couverture estimé rapporte 12.2~points de base par
mois, un ratio de Sharpe de 0.99, et le rendement croît avec la taille estimée du flux. La construction est \texttt{firm.fxflows}.

\section{Le fixing et ses flux}

Un fixing de référence (le livre~2, chapitre~17) fixe un taux de change de référence à partir des transactions réalisées dans une courte fenêtre de fixing. Les sociétés de gestion évaluent leurs portefeuilles et règlent
leurs couvertures à ce taux, si bien qu'un volume important d'ordres est exécuté au fixing par conception. Evans a étudié 21 devises sur une décennie. Autour du fixing WMR de 16~h,
les variations de prix présentaient une volatilité extraordinaire et une corrélation sérielle négative~: les prix poussés pendant la fenêtre avaient tendance à revenir après elle.

\begin{definition}[Flux du fixing]\label{def:s2:fx-flow-strategies:flow}
Un \emph{flux du fixing}\index{flux du fixing} est l'ordre net d'achat ou de vente d'une devise au fixing de référence, issu de mandats plutôt que d'opinions~: rééquilibrage des couvertures de change,
rééquilibrage d'indices et de portefeuilles, et ordres de clients exécutés au fixing~; sa taille et son sens peuvent souvent être estimés à l'avance à partir
d'informations publiques.
\end{definition}

Le fixing a aussi été le théâtre de pratiques abusives. Le 12 novembre 2014, la Financial Conduct Authority a infligé à cinq banques une amende totale de 1\,114\,918\,000~\pounds pour
des manquements sur le marché de change au comptant du G10 entre 2008 et 2013. Leurs traders avaient échangé des informations confidentielles sur les ordres de leurs clients et tenté de manipuler
les taux de fixing, y compris les fixings WM/Reuters et BCE de 16~h, et de déclencher les ordres stop-loss des clients. Estimer un flux à partir de données publiques et prendre position
par anticipation sur ce flux est légal. Négocier en s'appuyant sur la connaissance de l'ordre d'un client, ou s'entendre pour déplacer le fixing, ne l'est pas (chapitre~29).

\section{Rééquilibrage de fin de mois}

\begin{definition}[Estimation des flux de rebalancement]\label{def:s2:fx-flow-strategies:estimate}
Une \emph{estimation des flux de rebalancement}\index{estimation des flux de rebalancement} est une prévision de la devise que des investisseurs internationaux qui se couvrent doivent acheter ou
vendre à une date de rééquilibrage, à partir des rendements d'actifs du mois, des estimations de leurs positions et de leurs ratios de couverture.
\end{definition}

Le mécanisme est purement arithmétique. Un investisseur européen qui détient des actions américaines et couvre la moitié de leur valeur en dollars est vendeur de dollars à terme. Si les actions américaines
progressent de 5~\% au cours du mois, la position vaut davantage et la couverture couvre moins de la moitié. En fin de mois, l'investisseur vend davantage de dollars à terme pour rétablir
le ratio. Agrégé sur l'ensemble des investisseurs, un marché américain haussier signifie des ventes de dollars au fixing de fin de mois. Un investisseur américain qui couvre des actions étrangères fait l'inverse.

Les données publiques font apparaître ce schéma. Construire un indice du dollar à partir des taux quotidiens de midi à New York de six devises (euro, livre sterling, yen, franc suisse, dollars australien et
canadien), et prendre le rendement du S\&P 500 du début de chaque mois jusqu'à trois jours avant sa fin. Sur 332 mois, de 1999 à 2026, la
variation du dollar sur les deux derniers jours est corrélée à ce rendement à $-0.18$ ($t$ = $-3.3$). Vendre le dollar sur ces deux jours après un mois haussier
et l'acheter après un mois baissier, a rapporté 11.4~points de base par mois ($t$ = 3.4)
(la \cref{fig:s2:fx-flow-strategies:real}). Sur les trois derniers jours, le trade rapporte 7.4 ($t$ = 1.9), et sur les cinq derniers 13.7 ($t$ = 2.8). Les taux
sont relevés à midi à New York, plutôt qu'au fixing de Londres, si bien que le test n'en saisit le schéma qu'indirectement.

\begin{omfigure}
\begin{tikzpicture}
\begin{axis}[width=8cm,height=5.2cm,ybar,bar width=22pt,xlabel={\small $k$ derniers jours de négociation du mois},ylabel={\small trade (pb par mois)},
  tick label style={font=\footnotesize},xmin=1,xmax=6,ymin=0,ymax=16,xtick={2,3,5},grid=major,grid style={gray!25},table/col sep=comma,
  nodes near coords,every node near coord/.append style={font=\scriptsize}]
\addplot[fill=omThm!70,draw=omThm] table[x=k,y=bp]{figdata/strategies-2/15-fx-flow-strategies/real.csv};
\end{axis}
\end{tikzpicture}

\omcaption{Un trade sur le flux de couverture de fin de mois, sur données publiques~: vendre un indice du dollar à six devises pendant les $k$ derniers jours du mois lorsque le S\&P 500 a progressé
depuis le début du mois, l'acheter lorsque le S\&P a reculé~; rendement moyen par mois, de janvier 1999 à septembre 2026 (332 mois). D'après les taux H.10 de FRED et
de l'historique du S\&P 500 de Cboe~; statistiques dérivées seulement. Données~: \texttt{s2\_fetch\_monthend}.}\label{fig:s2:fx-flow-strategies:real}
\end{omfigure}

\texttt{firm.fxflows} construit explicitement le mécanisme (le \cref{lst:s2:fx-flow-strategies:model}). Sur 360 mois synthétiques, les investisseurs étrangers détiennent des actions nationales
couvertes selon un ratio de 0.5, et les investisseurs nationaux détiennent des actions étrangères couvertes à 0.3. Le flux de fin de mois déplace le taux de change dans les derniers jours
avant le fixing de 12.2~points de base en écart type, et la moitié du mouvement s'inverse ensuite. Le trader estime le flux avec des erreurs de 30~\% sur les
positions et les ratios de couverture~; l'estimation reste corrélée à la vraie valeur à 0.94. Suivre le sens du flux estimé a rapporté 12.2~points de base par
mois après frais ($t$ = 5.4), un ratio de Sharpe de 0.99. Prendre le contre-pied du mouvement dans les jours suivant la fin du mois a rapporté 10.2 ($t$ = 4.9). Par quintile de la taille du flux estimé,
le trade a rapporté 0.2, 12.0, 17.7, 9.5 et 21.8~points de base, contre des poussées attendues de 1.3, 4.5, 7.7, 12.8 et 23.4
(la \cref{fig:s2:fx-flow-strategies:quintiles}).

\begin{omfigure}
\begin{tikzpicture}
\begin{axis}[width=10cm,height=5.4cm,ybar,bar width=10pt,xlabel={\small quintile des mois selon la taille du flux estimé},ylabel={\small pb par mois},
  tick label style={font=\footnotesize},xmin=0.4,xmax=5.6,ymin=0,ymax=26,xtick={1,2,3,4,5},grid=major,grid style={gray!25},table/col sep=comma,
  legend style={font=\scriptsize,at={(0.5,-0.3)},anchor=north},legend columns=2]
\addplot[fill=gray!40,draw=gray] table[x=q,y=est]{figdata/strategies-2/15-fx-flow-strategies/quintiles.csv};
\addplot[fill=omThm!70,draw=omThm] table[x=q,y=realised]{figdata/strategies-2/15-fx-flow-strategies/quintiles.csv};
\legend{poussée attendue par l'estimation,rendement moyen du trade}
\end{axis}
\end{tikzpicture}

\omcaption{Le trade synthétique sur flux de fin de mois, par quintile de la taille du flux estimé~: la poussée que l'estimation implique et le rendement moyen du trade après
frais, 360 mois. Données~: \texttt{s2\_fxflows.results}.}\label{fig:s2:fx-flow-strategies:quintiles}
\end{omfigure}

\section{Comportement des banques centrales}

\begin{definition}[Trade de réaction aux banques centrales]\label{def:s2:fx-flow-strategies:cb}
Un \emph{trade de réaction aux banques centrales}\index{trade de réaction aux banques centrales} prend position sur la réponse prévisible d'une banque centrale à sa devise, par exemple
une intervention près des niveaux qu'elle a déjà défendus ou un rééquilibrage des réserves après de forts mouvements, et prend le contre-pied ou suit le mouvement selon le
bilan de la banque et sa capacité à agir.
\end{definition}

Les banques centrales sont les plus gros traders de flux qui existent, et certaines sont prévisibles. Une banque qui est intervenue pour stopper une dépréciation au-delà d'un certain niveau
peut recommencer. Un gestionnaire de réserves qui rééquilibre ses parts de devises après un mouvement important crée des flux comparables à ceux d'un hedger. Les trades sont rares et
la contrepartie peut imprimer de la monnaie. Un trader qui prend le contre-pied de la ligne affichée par une banque centrale parie sur sa volonté d'agir et sur ses ressources pour le faire. Le chapitre
ne donne aucun chiffre pour ce trade~: les relevés d'intervention sont publiés tardivement et sous des formes variées, et aucun n'est utilisé ici.

\section{Négocier des flux prévisibles}

Dans la fenêtre de fixing elle-même, le flux synthétique pousse le taux de 30~\% de son mouvement de fin de mois, et 40~\% de cette poussée s'inverse dans la demi-heure suivante.
Un fournisseur de liquidité qui prend le côté opposé de la poussée a rapporté 0.75~points de base par mois après un coût de 0.5~point de base ($t$ = 2.7), un ratio de Sharpe de
0.50. C'est la corrélation sérielle négative relevée par Evans, vue du côté qui fournit la liquidité plutôt que de celui qui en demande.

Trois conditions rendent un trade sur flux légitime et durable. Le flux doit être estimé à partir d'informations publiques. Il doit être grand par rapport à la
profondeur du marché sur la fenêtre. Et le trader doit accepter d'avoir tort les mois où l'estimation se trompe, car elle est bruitée et
d'autres traders peuvent se trouver du même côté.

\section{Fiches de stratégie}

\begin{strategyfile}{Flux de rééquilibrage de couverture de fin de mois}\label{strat:s2:fx-flow-strategies:monthend}
\sfield{Qui paie le trader, et pourquoi} Les investisseurs internationaux en actions et en obligations qui rééquilibrent leurs couvertures de change selon un calendrier fixe, quel que soit
le prix.
\sfield{Instruments et plateformes} Comptant et change à terme sur les principales paires, jusqu'au fixing de fin de mois.
\sfield{Signal} Les rendements du mois, par pays, pour les actions (et les obligations), multipliés par les positions estimées et les ratios de couverture.
\sfield{Dimensionnement et exécution} Proportionnel au flux estimé~; sortie au fixing ou après celui-ci.
\sfield{Coûts} Fourchettes proches du fixing, lorsqu'elles s'élargissent.
\sfield{Comment elle s'éteint} Les hedgers déplacent leur rééquilibrage hors du fixing~; trop de traders anticipent le même flux.
\sfield{Horizon, capacité, infrastructure} Jours~; estimations des flux par pays et par classe d'actifs.
\sfield{Backtester honnêtement} Positions et ratios de couverture tels qu'ils pouvaient être connus alors~; prix au moment du fixing.
\sfield{Sources} Melvin et Prins (2015)~; test sur données publiques~: 11.4~points de base par mois sur les deux derniers jours, 1999--2026~; ce chapitre~: 12.2 sur le
marché synthétique.
\end{strategyfile}

\begin{strategyfile}{Fourniture de liquidité sur la fenêtre de fixing}\label{strat:s2:fx-flow-strategies:fix}
\sfield{Qui paie, et pourquoi} Des ordres qui doivent être exécutés au fixing et qui poussent le taux dans la fenêtre.
\sfield{Instruments et plateformes} Comptant sur la fenêtre de fixing.
\sfield{Signal} Le déséquilibre des ordres de la fenêtre, au fur et à mesure de sa formation.
\sfield{Dimensionnement et exécution} Prendre la contrepartie dans la fenêtre~; clôturer dans l'heure.
\sfield{Coûts} Fourchettes de la fenêtre.
\sfield{Comment elle s'éteint} La méthodologie du fixing change et étale les ordres sur une fenêtre plus longue.
\sfield{Horizon, capacité, infrastructure} Minutes~; accès au comptant à faible latence.
\sfield{Backtester honnêtement} Données de tick dans la fenêtre~; aucune connaissance des ordres des clients.
\sfield{Sources} Evans (2018) sur le retournement autour du fixing~; ce chapitre~: 0.75~points de base par mois.
\end{strategyfile}

\begin{strategyfile}{Repli après intervention de la banque centrale}\label{strat:s2:fx-flow-strategies:cb}
\sfield{Qui paie, et pourquoi} Des spéculateurs qui poussent une devise contre le niveau défendu par la banque centrale et qui doivent se replier lorsqu'elle intervient.
\sfield{Instruments et plateformes} Comptant et options sur la paire.
\sfield{Signal} Distance aux niveaux des interventions passées~; communiqués officiels.
\sfield{Dimensionnement et exécution} Petite taille, avec des options pour borner la perte si le niveau cède.
\sfield{Coûts} Primes des options.
\sfield{Comment elle s'éteint} La banque centrale abandonne le niveau.
\sfield{Horizon, capacité, infrastructure} Jours à semaines.
\sfield{Backtester honnêtement} Interventions telles que publiées, souvent des semaines plus tard.
\sfield{Sources} Aucun chiffre de performance n'a été vérifié.
\end{strategyfile}

\begin{strategyfile}{Flux de couverture de change selon la performance des actions}\label{strat:s2:fx-flow-strategies:equityhedge}
\sfield{Qui paie, et pourquoi} Comme pour le flux de fin de mois, mesuré à partir de la performance relative de deux marchés d'actions.
\sfield{Instruments et plateformes} La paire formée par les devises des deux marchés.
\sfield{Signal} Performance des actions nationales contre celle des actions étrangères sur le mois.
\sfield{Dimensionnement et exécution} Derniers jours avant la fin du mois.
\sfield{Coûts} Fourchettes.
\sfield{Comment elle s'éteint} Les ratios de couverture changent~; d'autres flux dominent.
\sfield{Horizon, capacité, infrastructure} Jours.
\sfield{Backtester honnêtement} Les clôtures des deux marchés avant la transaction.
\sfield{Sources} Melvin et Prins (2015)~; aucun chiffre de performance vérifié au-delà des tests de ce chapitre.
\end{strategyfile}

\section{Tutoriel~: seize heures}\label{tut:s2:fx-flow-strategies:tut}

\begin{tutorial}
\textbf{Objectif.} Simuler les flux de couverture de fin de mois, les estimer, prendre position en amont de ces flux et fournir de la liquidité au fixing~; tester le schéma sur des données publiques.
\textbf{État final~:} les chiffres du texte et les deux figures.
\begin{tutsteps}
\item \textbf{Flux et transactions}.
\omcode{code/firm/fxflows/firm_fxflows.py}{46}{76}{Flux de fin de mois issus des rendements boursiers et des ratios de couverture, et les deux transactions.}\label{lst:s2:fx-flow-strategies:model}
\item \textbf{Résultats}.
\omcode{code/strategies-2/15-fx-flow-strategies/python/s2_fxflows.py}{39}{51}{La transaction sur le flux selon la taille estimée, le repli et la transaction sur le fixing.}\label{lst:s2:fx-flow-strategies:results}
\item \textbf{Exécuter} \texttt{s2\_fetch\_monthend.py} une fois, puis \texttt{results()} et \texttt{fig\_fxflows.py}.
\end{tutsteps}
\textbf{Pour aller plus loin.} Laisser les ratios de couverture dériver au fil des années~; ajouter un second flux, non corrélé (rééquilibrage d'indices), qui compense parfois le
premier~; ne négocier que les deux cinquièmes dont le flux estimé est le plus important.
\end{tutorial}

\section{Construction~: flux de change}\label{bld:s2:fx-flow-strategies:fxflows}

\begin{build}
\bfield{Rôle} Flux de couverture de fin de mois issus de la performance des actifs et des ratios de couverture, poussée et retournement sur la fenêtre de fixing, et les transactions qui les entourent.
\bfield{Interface} \texttt{FlowFXConfig(\dots)}, \texttt{simulate\_months(cfg)}, \texttt{flow\_trade(sim, cfg, scale)},
\texttt{fix\_liquidity(sim, cfg)}.
\bfield{Règles} Le signe du flux suit les couvertures~; l'estimation utilise des positions bruitées et des ratios de couverture~; coûts par transaction.
\bfield{Tests d'acceptation} \texttt{code/firm/fxflows/tests/}~: le signe et la taille du flux vérifiés à la main~; sans bruit, les transactions rapportent la poussée, ou son
retournement, nets des coûts.
\bfield{Pour aller plus loin} Plusieurs devises~; couvertures obligataires~; rééquilibrage d'indices.
\end{build}

\begin{omsources}
\item M.~Melvin et J.~Prins, «~Equity hedging and exchange rates at the London 4 p.m.\ fix~», \emph{Journal of Financial Markets} 22, 2015.
\item M.~D.~D.~Evans, «~Forex trading and the WMR Fix~», \emph{Journal of Banking and Finance} 87, 2018.
\item Financial Conduct Authority, communiqué de presse du 12 novembre 2014 sur les manquements de cinq banques en matière de change.
\item Board of Governors H.10 exchange rates, via FRED~; Cboe S\&P 500 index history.
\end{omsources}

\section{Exercices}

\begin{exercise}[$\star$]\label{exo:s2:fx-flow-strategies:1}
Un investisseur européen détient 100~millions de dollars d'actions américaines couvertes à 50~\%. Les actions américaines progressent de 5~\% sur le mois. Combien de dollars doit-il vendre à terme
à la fin du mois~?
\end{exercise}

\begin{exercise}[$\star$]\label{exo:s2:fx-flow-strategies:2}
Une transaction rapporte 11.4 points de base par mois avec $t$ = 3.4 sur 332 mois. Quel est son ratio de Sharpe annuel~?
\end{exercise}

\begin{exercise}[$\star$]\label{exo:s2:fx-flow-strategies:3}
Quel était le montant moyen des amendes infligées par la FCA par banque en novembre 2014~?
\end{exercise}

\begin{exercise}[$\star\star$]\label{exo:s2:fx-flow-strategies:4}
Pourquoi une hausse du marché américain des actions conduit-elle à des ventes de dollars en fin de mois~?
\end{exercise}

\begin{exercise}[$\star\star$]\label{exo:s2:fx-flow-strategies:5}
Pourquoi le test sur données publiques n'est-il qu'une approximation de la transaction sur le fixing~?
\end{exercise}

\begin{exercise}[$\star\star$]\label{exo:s2:fx-flow-strategies:6}
Qu'est-ce qui distingue une transaction légitime sur flux du comportement sanctionné par la FCA~?
\end{exercise}

\begin{exercise}[$\star\star\star$]\label{exo:s2:fx-flow-strategies:7}
\emph{Programmation.} Relancer avec \texttt{FlowFXConfig(est\_error=1.0)}. Comment évoluent la corrélation de l'estimation avec le flux réel et le rendement de la transaction~?
\end{exercise}

\begin{exercise}[$\star\star\star$]\label{exo:s2:fx-flow-strategies:8}
\emph{Trouver la faille.} «~La transaction sur flux fonctionne dans chaque cinquième de taille estimée, il faudrait donc la dimensionner de façon égale chaque mois.~»
\end{exercise}

\section{Problème~: seize heures}

\begin{problem}[{Problème du week-end --- flux de change prévisibles}]\label{pb:s2:fx-flow-strategies:1}
Les flux synthétiques du chapitre, les données publiques et les documents publics.

\medskip\noindent\textbf{Partie~I --- Le fixing.}
\begin{enumerate}
\item Qu'est-ce qu'un fixing de référence, et pourquoi y négocie-t-on tant d'ordres~?
\item Qu'a mis en évidence Evans autour du fixing~?
\item Qu'a constaté la FCA en 2014~?
\item Définir un \omterm{def:s2:fx-flow-strategies:flow}{flux du fixing}.
\end{enumerate}

\medskip\noindent\textbf{Partie~II --- Fin de mois.}
\begin{enumerate}[resume]
\item Expliquer le flux de rééquilibrage de la couverture.
\item Qu'ont trouvé Melvin et Prins~?
\item Donner les résultats du test sur données publiques.
\item Définir une \omterm{def:s2:fx-flow-strategies:estimate}{estimation des flux de rebalancement}.
\end{enumerate}

\medskip\noindent\textbf{Partie~III --- Le marché synthétique.}
\begin{enumerate}[resume]
\item Décrire les flux synthétiques et l'estimation du trader.
\item Donner le rendement de la transaction sur flux et le ratio de Sharpe.
\item Comment le rendement varie-t-il avec la taille estimée~?
\item Combien le fournisseur de liquidité gagne-t-il dans la fenêtre de fixing~?
\end{enumerate}

\medskip\noindent\textbf{Partie~IV --- Le verdict.}
\begin{enumerate}[resume]
\item Énoncer le \emph{résultat nommé}~: le rendement de la transaction sur flux de fin de mois en fonction de la taille de flux estimée.
\item Que rapporte le fading après la fin de mois~?
\item Définir un \omterm{def:s2:fx-flow-strategies:cb}{trade de réaction aux banques centrales} et son risque.
\item Qu'est-ce qui rend une transaction sur flux légitime~?
\item Comment backtester honnêtement la transaction de fin de mois~?
\item Quelle fiche de stratégie exige l'exécution la plus rapide~?
\item Quel lien ce chapitre entretient-il avec les flux de rebalancement du chapitre~5~?
\item En une phrase~: qui paie le trader sur flux~?
\end{enumerate}
\end{problem}

\section{Questions d'entretien}

\begin{interviewq}[$\star$ \iqroles{trader}]\label{iq:s2:fx-flow-strategies:1}
Qu'est-ce que le fixing WM/Reuters de 16~h, et pourquoi a-t-il de l'importance~?
\end{interviewq}

\begin{interviewq}[$\star\star$ \iqroles{chercheur}]\label{iq:s2:fx-flow-strategies:2}
Comment estimeriez-vous les flux de couverture de change de fin de mois~?
\end{interviewq}

\begin{interviewq}[$\star\star$ \iqroles{trader}]\label{iq:s2:fx-flow-strategies:3}
Un client vous demande d'acheter 2~milliards de dollars au fixing. Que pouvez-vous faire avant le fixing, et que vous est-il interdit de faire~?
\end{interviewq}

\begin{interviewq}[$\star\star$ \iqroles{risque}]\label{iq:s2:fx-flow-strategies:4}
Quels sont les risques d'un portefeuille de flux de fin de mois~?
\end{interviewq}

\begin{interviewq}[$\star\star$ \iqroles{développeur}]\label{iq:s2:fx-flow-strategies:5}
Concevez un outil de surveillance des transactions autour des fixings de référence qu'une équipe de conformité accepterait.
\end{interviewq}

\begin{interviewq}[$\star\star\star$ \iqroles{chercheur}]\label{iq:s2:fx-flow-strategies:6}
Avec des positions $H$, un ratio de couverture $h$ et un rendement d'actif $R$, montrer que rétablir la couverture exige une transaction de change de $hHR$, et trouver l'erreur de l'estimation
lorsque $H$ et $h$ sont chacun décalés d'un facteur $(1 + \epsilon)$.
\end{interviewq}
